\documentclass[11pt, letterpaper]{article}

% ── Packages ──
\usepackage[utf8]{inputenc}
\usepackage[T1]{fontenc}
\usepackage{amsmath, amssymb, amsthm, mathtools}
\usepackage{microtype}
% Body + math: Computer Modern (LaTeX/AMS default) via Latin Modern
\usepackage{lmodern}
% Sans-serif (used for section headings): Latin Modern Sans
\renewcommand{\sfdefault}{lmss}
\usepackage{geometry}
\usepackage{hyperref}
\usepackage{enumitem}
\usepackage{booktabs}
\usepackage{multirow}
\usepackage{array}
\usepackage{caption}
\usepackage[svgnames]{xcolor}
\usepackage{titling}
\usepackage{titlesec}

\geometry{margin=1.05in}
\hypersetup{
  colorlinks=true,
  linkcolor=MidnightBlue,
  citecolor=MidnightBlue,
  urlcolor=MidnightBlue,
  pdftitle={A New Lower Bound for the Supremum of Autoconvolutions},
  pdfauthor={Andrei Piterbarg, Jai Bajaj, Derrick Vincent},
  pdfsubject={A computer-assisted lower bound for the autoconvolution constant}
}
\captionsetup{font=small, labelfont=bf}
\setlist{topsep=4pt, itemsep=2pt, parsep=0pt}
\setlength{\parindent}{1.1em}
\setlength{\parskip}{0pt}
\allowdisplaybreaks
\setlength{\droptitle}{-2.8em}
\pretitle{\begin{center}\LARGE}
\posttitle{\par\end{center}\vspace{-0.65em}}
\preauthor{\begin{center}\normalsize}
\postauthor{\par\end{center}\vspace{-0.45em}}
\predate{\begin{center}\small}
\postdate{\par\end{center}\vspace{0.6em}}

% Sections: centered, serif bold (matches the Abstract title font)
\titleformat{\section}[block]
  {\normalfont\Large\bfseries\centering}
  {\thesection.}{0.55em}{}
% Subsections: left-aligned, serif bold (matches the Abstract title font)
\titleformat{\subsection}[hang]
  {\normalfont\large\bfseries}
  {\thesubsection.}{0.5em}{}
\titlespacing*{\section}{0pt}{2.6ex plus 0.8ex minus 0.3ex}{1.4ex}
\titlespacing*{\subsection}{0pt}{1.8ex plus 0.5ex minus 0.2ex}{0.8ex}

% Abstract: match the \section heading style (centered, \Large\bfseries,
% serif) so the "Abstract" title is the same size as "1. Introduction".
% Body keeps the article-class \small + quotation framing.
\renewenvironment{abstract}{%
  \vspace{2.6ex plus 0.8ex minus 0.3ex}%
  \begin{center}{\normalfont\Large\bfseries\abstractname}\end{center}%
  \vspace{1.4ex}%
  \small\quotation
}{%
  \endquotation
}

% ── Theorem environments ──
\newtheoremstyle{paperplain}
  {8pt}{8pt}{\itshape}{}{\bfseries}{.}{0.5em}{}
\newtheoremstyle{paperdef}
  {8pt}{8pt}{\normalfont}{}{\bfseries}{.}{0.5em}{}
\theoremstyle{paperplain}
\newtheorem{theorem}{Theorem}[section]
\newtheorem{lemma}[theorem]{Lemma}
\newtheorem{proposition}[theorem]{Proposition}
\newtheorem{corollary}[theorem]{Corollary}
\newtheorem{claim}[theorem]{Claim}
\theoremstyle{paperdef}
\newtheorem{definition}[theorem]{Definition}
\newtheorem{remark}[theorem]{Remark}
\newtheorem{example}[theorem]{Example}

% ── Macros ──
\newcommand{\R}{\mathbb{R}}
\newcommand{\N}{\mathbb{N}}
\newcommand{\Z}{\mathbb{Z}}
\newcommand{\Q}{\mathbb{Q}}
\newcommand{\Linf}{L^\infty}
\newcommand{\supp}{\operatorname{supp}}
\newcommand{\tv}{\operatorname{tv}}
\newcommand{\rev}{\operatorname{rev}}
\newcommand{\conv}{\operatorname{conv}}
\newcommand{\Karc}{K_{\mathrm{arc}}}
\newcommand{\Kms}{K_{\mathrm{ms}}}
\newcommand{\Mcert}{M_{\mathrm{cert}}}

% ── Hyphenation hints for long compound terms ──
\hyphenation{
  Matolcsi
  Vinuesa
  Bochner-ad-mis-si-ble
  arc-sine
  auto-con-vo-lu-tion
  auto-cor-re-la-tion
  Cloninger
  Stein-er-berg-er
  inter-val
  rig-or-ous
  for-mal-i-za-tion
  multi-scale
  AlphaEvolve
}

% ══════════════════════════════════════════════════════════════════════
\title{\texorpdfstring{%
  A New Lower Bound for the\\
  Supremum of Autoconvolutions
}{A New Lower Bound for the Supremum of Autoconvolutions}}
\author{Andrei Piterbarg, Jai Bajaj and Derrick Vincent}
\date{}

\begin{document}
\maketitle

\begin{abstract}
We establish a lower bound on the
autoconvolution constant
$C_{1a}=\inf_f \|f*f\|_{\Linf}/(\int f)^2$, taken over nonnegative
$f\in L^1(\R)$ supported on $(-1/4,1/4)$:
\[
  C_{1a}\;\ge\;1292/1000\;=\;1.292.
\]
This improves on the previously announced bound of
$1.2802$ due to Cloninger and Steinerberger~\cite{CS17} and on the
rigorous analytic bound of $1.27481$ established by Matolcsi and
Vinuesa~\cite{MV10}.  The argument extends the Matolcsi--Vinuesa
master inequality: the single arcsine kernel is
replaced by a convex combination of three arcsine kernels, and the
trigonometric multiplier $G$ is re-optimized as a $200$-cosine
expansion.
\par\medskip
The five resulting functionals are evaluated rigorously in
\texttt{flint.arb} interval arithmetic at $256$-bit precision.  The
no-$z_1$ quadratic master inequality, applied to the rounded rational
anchors, yields strict failure at $M=1292/1000$ with margin
$\tau-\Phi(M)\ge 9.6\times 10^{-5}$.  The analytic reduction is
formalized in Lean~4: the headline theorem in
\texttt{lean/Sidon/MultiScale.lean} reaches a single user axiom
(the three-scale MV master inequality, stated at the rational slacks);
the quadratic inversion and the slack-soundness statements are Lean
theorems (see Section~\ref{sec:proof}).
\end{abstract}

% ======================================================================
\section{Introduction}\label{sec:intro}
% ======================================================================

\subsection{The Constant \texorpdfstring{$C_{1a}$}{C1a}}\label{subsec:intro-constant}

Let $\mathcal{F}=\{f\in L^1(\R) : f\ge 0,\
\supp(f)\subseteq(-1/4,1/4),\ \int f>0\}$.  The
\emph{autoconvolution constant} is
\[
  C_{1a}
  \;=\;
  \inf_{f\in\mathcal{F}} R(f),
  \qquad
  R(f) \;:=\; \frac{\|f*f\|_{\Linf}}{\left(\int f\right)^2}.
\]
By homogeneity we may normalize $\int f=1$, in which case
$R(f)=\|f*f\|_{\Linf}\ge 1$.  Heuristically, $R(f)$ measures the
peak height of the self-convolution $f*f$ relative to its $L^1$
mass, so $C_{1a}$ identifies the flattest autoconvolution achievable
in the class $\mathcal{F}$.  This constant is the continuous
analogue of the Sidon-set extremal problem~\cite{ET41,MO2009};
improving its lower bound tightens known asymptotic results on
$B_2^{+}[1]$ sets and related Sidon variants.

\subsection{Known Bounds}\label{subsec:intro-bounds}

Prior lower bounds on $C_{1a}$ are summarized in
Table~\ref{tab:published-bounds}.  Matolcsi--Vinuesa~\cite{MV10}
established the rigorous analytic value $1.27481$ in 2010, and
Cloninger and Steinerberger~\cite{CS17} announced a computational
improvement to $1.2802$ via a discrete branch-and-prune search.  All
subsequent progress on the constant has been on the upper-bound
side, where a recent AI-assisted search lowered the ceiling to
$1.5029$~\cite{AlphaEvolve25}.  The present work raises the
lower bound to $1.292$ (Theorem~\ref{thm:main}), which we refer to as
the \emph{Piterbarg--Bajaj--Vincent Bound} (henceforth, the
\emph{PBV~bound}), by refining the kernel class entering the
Matolcsi--Vinuesa dual framework.

\begin{table}[ht]
\centering
\begin{tabular}{lcl}
  \toprule
  \textbf{Bound} & \textbf{Value} & \textbf{Reference} \\
  \midrule
  Upper bound & $\le 1.5029$  & Georgiev et al.~\cite{AlphaEvolve25} (AlphaEvolve) \\
  Upper bound & $\le 1.50992$ & Matolcsi \& Vinuesa~\cite{MV10} \\
  Lower bound & $\ge 1.262$   & Martin \& O'Bryant~\cite{MO2009} \\
  Lower bound & $\ge 1.27481$ & Matolcsi \& Vinuesa~\cite{MV10} \\
  Lower bound & $\ge 1.2802$  & Cloninger \& Steinerberger~\cite{CS17} \\
  Lower bound & $\ge 1.292$   & \textbf{This work} \\
  \bottomrule
\end{tabular}
\caption[Published bounds on C1a in the normalization used here.]{Published bounds on $C_{1a}$ in the normalization used here.}
\label{tab:published-bounds}
\end{table}

\subsection{Strategy and Main Result}\label{subsec:intro-strategy}

The Matolcsi--Vinuesa~\cite{MV10} dual framework
(Section~\ref{sec:mv}) parametrizes a Bochner-admissible kernel $K$ and
a cosine multiplier $G$, producing a quadratic lower bound on $R(f)$
in five real functionals $(k_1,K_2,S_1,m_G,a)$, of which $(k_1,K_2)$
are extracted from the kernel and $(S_1,m_G,a)$ from the multiplier
(with $a=(4/u)\,m_G^{\,2}/S_1$ a derived gain functional).  MV~2010 takes $K$ to
be the arcsine kernel at half-width $\delta=138/1000$ and $G$ a
$119$-cosine expansion, achieving $R(f)\ge 1.27481$.  The present
construction (i)~replaces the single arcsine by a three-scale convex
combination $\Kms$ at half-widths $(138,55,25)/1000$ with weights
$(85,10,5)/100$; (ii)~re-optimizes $G$ as a $200$-cosine expansion
against $\Kms$; (iii)~certifies the five functionals in
\texttt{flint.arb} interval arithmetic; and (iv)~solves the no-$z_1$
quadratic master inequality with the certified rationals.  The
quadratic inequality $\Phi(M)\ge\tau$ fails strictly at the rational
witness $M=1292/1000$, with margin $9.6\times 10^{-5}$ from the
rational anchors, giving the main result.

\begin{theorem}[Main result]
\label{thm:main}
Let $f:\R\to\R$ be nonnegative with
$\supp(f)\subseteq(-1/4,1/4)$, $\int f>0$, and
$\|f*f\|_{\Linf}<\infty$.  Then
\[
  \frac{\|f*f\|_{\Linf}}{\left(\int f\right)^2}
  \;\ge\;
  \frac{1292}{1000}.
\]
In particular $C_{1a}\ge 1292/1000=1.292$.
\end{theorem}

The proof is assembled in Section~\ref{sec:proof}; all numerical
anchors are reproducible from the accompanying certificate, and a
Lean~4 mechanization is summarized at the end of
Section~\ref{sec:proof}.

% ======================================================================
\section{The Matolcsi--Vinuesa Master Inequality}\label{sec:mv}
% ======================================================================

We summarize the dual framework of~\cite[\S 3]{MV10}; proofs of the
four ingredient lemmas and of the resulting quadratic are recorded
in that reference.  Throughout, $f\in\mathcal{F}$ with $\int f=1$,
$\widehat{h}(\xi)=\int h(x)e^{-2\pi i\xi x}\,dx$, and
$M=R(f)=\|f*f\|_{\Linf}$.  The framework converts the desired lower
bound on $M$ into a quadratic inequality whose coefficients depend on
two auxiliary objects of our choice, a positive-definite kernel $K$
and a trigonometric multiplier $G$, so that any rigorous numerical
control over $(K,G)$ translates into a rigorous lower bound on $M$.

\subsection{Admissible Kernels and Multipliers}\label{subsec:mv-defs}

\begin{definition}[Bochner-admissible kernel]
\label{def:bochner}
Fix $\delta\in(0,1/4]$ and set $u=1/2+\delta$.  A function
$K:\R\to\R$ is \emph{Bochner-admissible at scale $\delta$} if
\begin{enumerate}[label=\textnormal{(K\arabic*)}, leftmargin=2.4em]
\item $K\ge 0$ a.e.\ on $\R$;
\item $K(-x)=K(x)$ for every $x\in\R$;
\item $\supp(K)\subseteq[-\delta,\delta]$ and $\int K=1$;
\item the periodic Fourier coefficients
      $\widetilde{K}(j):=\widehat{K}(j/u)$ satisfy
      $\widetilde{K}(j)\ge 0$ for every $j\in\Z$.
\end{enumerate}
We write $K_2:=\|K\|_{L^2(\R)}^2$.
\end{definition}

\begin{definition}[Admissible cosine multiplier]
\label{def:multiplier}
Fix $u$ and $N\in\N$.  A trigonometric polynomial
$G(x)=\sum_{j=1}^{N}a_j\cos(2\pi jx/u)$ with $a_j\in\R$ is
\emph{admissible} if there is $m_G>0$ with $G(x)\ge m_G$ for every
$x\in[0,1/4]$.  Define
\begin{equation}\label{eq:S1-a}
  S_1 \;=\; \sum_{j=1}^{N}\frac{a_j^{\,2}}{\widetilde{K}(j)},
  \qquad
  a \;=\; \frac{4}{u}\cdot\frac{m_G^{\,2}}{S_1}.
\end{equation}
\end{definition}

\subsection{The Master Inequality and Its Inversion}\label{subsec:mv-master}

\begin{theorem}[Master inequality, $z_1$-free form]
\label{thm:mv-master}
Let $K$ be Bochner-admissible at scale $\delta$, $u=1/2+\delta$,
and $G$ admissible of degree $N$ with constant $m_G$, defining $S_1$
and $a$ by~\eqref{eq:S1-a}.  Then for every $f\in\mathcal{F}$ with
$\int f>0$,
\begin{equation}\label{eq:master}
  M \;+\; 1 \;+\; \sqrt{(M-1)(K_2-1)}
  \;\ge\;
  \frac{2}{u} \;+\; a.
\end{equation}
\end{theorem}

\begin{proof}
Matolcsi--Vinuesa~\cite[Eq.~(7), Lemma~3.1]{MV10} prove the sharper
form
\begin{equation}\label{eq:mv-sharp}
  M + 1 + 2 z_1 k_1
  + \sqrt{(M-1-2z_1^{\,2})(K_2-1-2k_1^{\,2})}
  \;\ge\;
  \frac{2}{u} + a,
\end{equation}
where $k_1=\widetilde{K}(1)$ and $z_1=\widehat{f}(1/u)$.  The
\emph{$z_1$-free} designation reflects that $z_1$ is a Fourier
coefficient of the (unknown) extremal $f$, so it cannot be controlled
analytically; the Cauchy--Schwarz step below absorbs $z_1$ (together
with $k_1$) into the kernel-side quantity $K_2$, leaving an inequality
that depends only on quantities we can certify.  Their proof
uses the arcsine kernel only through Bochner-admissibility
(Definition~\ref{def:bochner}); the same derivation therefore applies
to any admissible $K$, in particular to the convex
combination~$\Kms$.  The radicand of~\eqref{eq:mv-sharp} is nonnegative
by the same Parseval-type bounds as in~\cite[\S 3]{MV10}: Bessel's
inequality gives $2z_1^{\,2}\le M-1$ for $f\in\mathcal{F}$ with
$\int f=1$, and $2k_1^{\,2}\le K_2-1$ holds for any Bochner-admissible
$K$ with $\int K=1$.  Inequality~\eqref{eq:master} then follows from
\eqref{eq:mv-sharp} by the Cauchy--Schwarz estimate
$\sqrt{ac}+\sqrt{bd}\le\sqrt{(a+b)(c+d)}$ applied to the nonnegative
quadruple
$(a,b,c,d)=(2z_1^{\,2},\,M-1-2z_1^{\,2},\,2k_1^{\,2},\,K_2-1-2k_1^{\,2})$.
\end{proof}

\begin{lemma}[Quadratic inversion]
\label{lem:inversion}
For fixed $K_2\ge 1$, the function
\[
  \Phi(M)\;:=\;M+1+\sqrt{(M-1)(K_2-1)}
\]
is continuous and strictly increasing on $[1,\infty)$ with
$\Phi(1)=2$.  For any $\tau\ge 2$ the equation $\Phi(M)=\tau$ has a
unique solution $M_*\in[1,\infty)$, and every $f\in\mathcal{F}$ with
$\Phi(R(f))\ge\tau$ satisfies $R(f)\ge M_*$.  In particular, with
$\tau=2/u+a$, $C_{1a}\ge M_*$.
\end{lemma}

\begin{proof}
For $M>1$, $\Phi'(M)=1+\sqrt{(K_2-1)/(M-1)}/2>1$, so $\Phi$ is
strictly increasing on $(1,\infty)$ and continuous at $M=1$ from the
right.  Existence, uniqueness, and the conclusion follow from the
intermediate value theorem and Theorem~\ref{thm:mv-master}.
\end{proof}

The MV~2010 single-arcsine optimum yields $M_*\approx 1.2748$.  The
remainder of the paper raises $M_*$ by enlarging the kernel class.

% ======================================================================
\section{The Three-Scale Arcsine Kernel}\label{sec:kernel}
% ======================================================================

\subsection{Single-Scale Arcsine and the Sonine Identity}\label{subsec:kernel-single}

For $\delta\in(0,1/4]$, let
\[
  \eta_\delta(x)
  \;:=\;
  \frac{1}{\delta}\cdot
  \frac{2/\pi}{\sqrt{1-(2x/\delta)^2}}\,
  \mathbf{1}_{|x|<\delta/2}
\]
denote the arcsine density on $(-\delta/2,\delta/2)$ (the
pushforward of the arcsine density on $(-1/2,1/2)$ under
$u\mapsto\delta u$).  By the Sonine--Bessel
identity~\cite[\S 13.46]{Watson}, $\widehat{\eta_\delta}(\xi)
=J_0(\pi\delta\xi)$.  The \emph{arcsine kernel at scale $\delta$} is the
auto-convolution
\[
  \Karc(\delta;\cdot)
  \;:=\;
  \eta_\delta*\eta_\delta,
\]
which is nonnegative and even, supported on $[-\delta,\delta]$, with
$\int\Karc(\delta;\cdot)=(\int\eta_\delta)^2=1$.  The convolution
theorem then gives
\begin{equation}\label{eq:sonine}
  \widehat{\Karc(\delta;\cdot)}(\xi)
  \;=\;
  \widehat{\eta_\delta}(\xi)^{2}
  \;=\; J_0(\pi\delta\xi)^2
  \;\ge\; 0.
\end{equation}
In particular each $\Karc(\delta;\cdot)$ is Bochner-admissible at
scale $\delta$, with periodic coefficients
$\widetilde{\Karc}(\delta;j)=J_0(\pi j\delta/u)^2$ for $u\ge 2\delta$.

\subsection{The Convex Combination}\label{subsec:kernel-combo}

\begin{definition}[Three-scale kernel]
\label{def:Kms}
Fix
\begin{equation}\label{eq:anchors}
  (\delta_1,\delta_2,\delta_3)
  \;=\;
  \bigl(\tfrac{138}{1000},\tfrac{55}{1000},\tfrac{25}{1000}\bigr),
  \qquad
  (\lambda_1,\lambda_2,\lambda_3)
  \;=\;
  \bigl(\tfrac{85}{100},\tfrac{10}{100},\tfrac{5}{100}\bigr),
\end{equation}
and set $u=1/2+\delta_1=638/1000$.  The
\emph{three-scale arcsine kernel} is
\[
  \Kms(x)
  \;:=\;
  \sum_{i=1}^{3}\lambda_i\,\Karc(\delta_i;x),
  \qquad x\in\R.
\]
\end{definition}

\begin{table}[ht]
\centering
\begin{tabular}{ccccc}
  \toprule
  $i$ & $\delta_i$ & $\lambda_i$ & $u$ & $N$ \\
  \midrule
  $1$ & $138/1000$ & $85/100$ & \multirow{3}{*}{$638/1000$} & \multirow{3}{*}{$200$} \\
  $2$ & $55/1000$  & $10/100$ & & \\
  $3$ & $25/1000$  & $5/100$  & & \\
  \bottomrule
\end{tabular}
\caption[Parameters of the three-scale arcsine kernel and the cosine multiplier.]{Parameters of the three-scale arcsine kernel and the cosine
multiplier.  The period $u$ and the cosine degree $N$ are shared
across the three scales.}
\label{tab:kernel-params}
\end{table}

\subsection{Admissibility}\label{subsec:kernel-admissibility}

\begin{theorem}[Admissibility]
\label{thm:admissibility}
The kernel $\Kms$ is Bochner-admissible at scale
$\delta_1=138/1000<1/4$, with periodic coefficients
\begin{equation}\label{eq:Ktilde}
  \widetilde{\Kms}(j)
  \;=\;
  \sum_{i=1}^{3}\lambda_i\,
       J_0\!\left(\frac{\pi j\delta_i}{u}\right)^{\!2}
  \;\ge\; 0,
  \qquad j\in\Z.
\end{equation}
\end{theorem}

\begin{proof}
We check conditions (K1)--(K4) of Definition~\ref{def:bochner}.
\begin{itemize}[leftmargin=2em]
\item \emph{(K1)~Nonnegativity, (K2)~evenness.}  Each
  $\lambda_i\Karc(\delta_i;\cdot)\ge 0$ and is even; nonnegative sums
  of even nonnegative functions are even and nonnegative.
\item \emph{(K3)~Support and unit mass.}
  $\supp\Karc(\delta_i;\cdot)\subseteq
  [-\delta_i,\delta_i]\subseteq[-\delta_1,\delta_1]$ for $i\ge 1$, so
  $\supp\Kms\subseteq[-\delta_1,\delta_1]$; and
  $\int\Kms=\sum_i\lambda_i\int\Karc(\delta_i)=\sum_i\lambda_i=1$.
\item \emph{(K4)~Fourier positivity.}  Linearity of the Fourier
  transform and~\eqref{eq:sonine} give
  \[
    \widehat{\Kms}(\xi)
    \;=\;
    \sum_i\lambda_i J_0(\pi\delta_i\xi)^2 \;\ge\; 0,
  \]
  hence~\eqref{eq:Ktilde}.\qedhere
\end{itemize}
\end{proof}

The motivation is structural: in~\eqref{eq:S1-a}, $S_1$ blows up at
frequencies $j$ where the single-scale $J_0(\pi j\delta/u)^2$ is close
to a Bessel zero.  Mixing three scales shifts the zeros of each
summand, so the sum stays bounded away from zero on $1\le j\le 200$
(numerically $\min_j\widetilde{\Kms}(j)\ge 2.08\times 10^{-4}$).  This
drops $S_1$ from $\approx 87.4$ (single-scale) to $\le 29.841$
(Lemma~\ref{lem:S1}), and proportionally raises the gain~$a$.

% ======================================================================
\section{The Rigorous Numerical Certificate}\label{sec:cert}
% ======================================================================

\subsection{The Five Certified Anchors}\label{subsec:cert-anchors}

For the master inequality applied to $(\Kms,G)$ we need rigorous
bounds on five real functionals.  We replace each by a rational
\emph{anchor}: a floor when the functional enters the master
inequality on the side that weakens the bound under understatement
(i.e.\ with a $\ge$ sign), a ceiling on the side that weakens it
under overstatement ($\le$), in both cases obtained by outward
rounding of the certifier's arb interval enclosure.  Anchoring at
exact rationals lets the final inversion (Section~\ref{sec:proof}) be
carried out in exact rational arithmetic, independent of the
floating-point pipeline that produced the enclosures.  The cosine multiplier
$G(x)=\sum_{j=1}^{N}a_j\cos(2\pi jx/u)$ with $N=200$ has rational
coefficients $a_j\in\Q$ with denominator $10^{8}$, obtained by
rounding the solution of the convex QP that minimizes
$\sum_j a_j^{\,2}/\widetilde{\Kms}(j)$ subject to $G\ge 1$ on a
uniform grid of $5001$ points in $[0, 1/4]$ (solved by MOSEK).  The
five lemmas below are arb-interval computations~\cite{Flint}; all
certified outputs are rounded outward to exact rationals.  Only four
of the five anchors ($K_2$, $S_1$, $m_G$, $a$) enter the
$z_1$-free form~\eqref{eq:master} used in Section~\ref{sec:proof};
the kernel-mass moment $k_1$ is recorded because the certifier
internally uses the sharper inequality~\eqref{eq:mv-sharp} (which
sharpens $M_{\rm cert}$ from $1.292$ to $\approx 1.29232$ before
rounding to the rational headline target).

\begin{lemma}[Kernel-mass moment]
\label{lem:k1}
\(
  k_1 \;:=\; \widehat{\Kms}(1)
  \;=\; \sum_{i=1}^{3}\lambda_i J_0(\pi\delta_i)^2
  \;\ge\; 9212/10000.
\)
\end{lemma}

\begin{proof}
At $256$-bit precision, \texttt{arb.bessel\_j(0)} evaluated at exact
rational $\pi\delta_i$ returns balls of radius $<7\times 10^{-77}$;
the certified sum is $k_1\ge 0.92124658$, with slack
$\ge 4.6\times 10^{-5}$ above $9212/10000$.
\end{proof}

\begin{lemma}[Kernel-energy moment]
\label{lem:K2}
\(
  K_2 \;=\; \|\Kms\|_{L^2(\R)}^{2}
  \;\le\; 47897/10000.
\)
\end{lemma}

\begin{proof}
Split
$K_2 = 2\int_0^{T}\widehat{\Kms}(\xi)^2\,d\xi
+ 2\int_{T}^{\infty}\widehat{\Kms}(\xi)^2\,d\xi$ with $T=10^{5}$.
The bulk integral is computed by arb adaptive Gauss--Legendre
quadrature and certified at $4.78882342\ldots$ with arb radius
below $10^{-10}$.  For the tail, the standard uniform bound
$|J_0(z)|^2\le 2/(\pi z)$~\cite[\S 7.21]{Watson} holds for $z\ge 1$;
since $\xi\ge T=10^{5}$ in the tail forces $z=\pi\delta_i\xi\ge
\pi(25/1000)(10^{5})\approx 7854$, the bound applies, and applied to
both Bessel factors gives
$J_0(\pi\delta_i\xi)^2 J_0(\pi\delta_j\xi)^2
\le 4/(\pi^4\delta_i\delta_j\xi^2)$, whence
$\widehat{\Kms}(\xi)^2 = \sum_{i,j}\lambda_i\lambda_j
   J_0(\pi\delta_i\xi)^2 J_0(\pi\delta_j\xi)^2
\le 4C^2/(\pi^4\xi^2)$ with
$C=\sum_i\lambda_i/\delta_i\approx 9.978$.  Integrating
$\int_T^\infty \xi^{-2}\,d\xi = 1/T$ and doubling for both half-lines,
$2\int_T^\infty\widehat{\Kms}^2 \le 8C^2/(\pi^4 T) \le 8.19\times 10^{-5}$.
Summing,
$K_2\in[4.78882342,\,4.78890519]\subseteq[0,\,47897/10000]$, slack
$\ge 7.9\times 10^{-4}$.
\end{proof}

\begin{lemma}[Multiplier denominator]
\label{lem:S1}
\(
  S_1 \;=\;
  \sum_{j=1}^{200}\frac{a_j^{\,2}}{\widetilde{\Kms}(j)}
  \;\le\; 29841/1000.
\)
\end{lemma}

\begin{proof}
Each $\widetilde{\Kms}(j)$ is evaluated in arb from~\eqref{eq:Ktilde}
with radius $<10^{-70}$; dividing exact rational squares $a_j^{\,2}$
yields a certified sum $S_1\le 29.840907$, slack
$\ge 9.3\times 10^{-5}$.
\end{proof}

\begin{lemma}[Pointwise lower bound of $G$]
\label{lem:mG}
\(
  m_G \;:=\; \min_{x\in[0,1/4]}G(x) \;\ge\; 998/1000.
\)
\end{lemma}

\begin{proof}
Partition $[0,1/4]$ into $n=32768$ closed cells of half-width
$r=1/(8n)\approx 3.81\times 10^{-6}$.  On each cell centred at $c$,
the second-order Taylor expansion of $G$ yields the rigorous arb
enclosure
\[
  G(\mathrm{cell})\;\subseteq\;
    G(c)\;+\;G'(c)\cdot[-r,r]\;+\;G''(\mathrm{cell})\cdot[-r^{2}/2,\,r^{2}/2],
\]
with $G''$ evaluated as an arb interval over the entire cell from
\[
  G''(x) \;=\; -\sum_{j=1}^{200}a_j\,(2\pi j/u)^{2}\cos(2\pi jx/u).
\]
The minimum of the lower endpoints over all $n$ cells gives the
certified bound $m_G\ge 0.99997987$, slack $\ge 1.9\times 10^{-3}$
above $998/1000$.
\end{proof}

\begin{lemma}[Gain functional]
\label{lem:a}
\(
  a \;=\; (4/u)\cdot m_G^{\,2}/S_1 \;\ge\; 20925/100000.
\)
\end{lemma}

\begin{proof}
Substituting the rational anchors of
Lemmas~\ref{lem:S1} and~\ref{lem:mG} into the definition of $a$
in~\eqref{eq:S1-a},
\[
  a
  \;\ge\;
  \frac{4}{638/1000}\cdot\frac{(998/1000)^{2}}{29841/1000}
  \;=\;
  \frac{4000\cdot 998^{2}}{638\cdot 29841\cdot 10^{3}}
  \;=\;
  \frac{3984016000}{19038558000}
  \;=\; 0.20926\ldots,
\]
hence $a\ge 20925/100000$, slack $\ge 1.0\times 10^{-5}$.  (The
underlying arb interval evaluation returns the tighter enclosure
$a\ge 0.21009214$ when $m_G$ and $S_1$ are coupled rather than
rounded separately; we use only the looser rational floor in what
follows.)
\end{proof}

\begin{table}[ht]
\centering
\begin{tabular}{lcccc}
  \toprule
  Functional & Symbol & Direction & Rational bound & Decimal \\
  \midrule
  Kernel-mass moment   & $k_1$ & $\ge$ & $9212/10000$    & $0.9212$ \\
  Kernel-energy        & $K_2$ & $\le$ & $47897/10000$   & $4.7897$ \\
  Multiplier denominator & $S_1$ & $\le$ & $29841/1000$  & $29.841$ \\
  Multiplier minimum   & $m_G$ & $\ge$ & $998/1000$      & $0.998$  \\
  Gain                 & $a$   & $\ge$ & $20925/100000$  & $0.20925$ \\
  \bottomrule
\end{tabular}
\caption[Certified rational anchors for the three-scale kernel.]{Certified rational anchors for the three-scale kernel
$\Kms$ at parameters~\eqref{eq:anchors} with $u=638/1000$ and the
$200$-cosine multiplier $G$.}
\label{tab:anchors}
\end{table}

\subsection{Rescaling Invariance}\label{subsec:cert-rescaling}

Although the proof of Lemma~\ref{lem:mG}
produces $m_G\ge 998/1000$, the master inequality is invariant
under $G\mapsto G/m_G$: rescaling multiplies both $S_1$ and $m_G^{\,2}$
by $1/m_G^{\,2}$, leaving $a$ unchanged.  The certified rational
bound $m_G\ge 998/1000$ is therefore sufficient for the conclusion,
and the lemma is not sharpened further.

\subsection{Reproducibility}\label{subsec:cert-repro}

All five anchors are produced by the reproducible script and the
associated certificate JSON: each entry is the outward-rounded enclosure
returned by \texttt{flint.arb} at $256$-bit precision; full pipeline
inputs (the rational $a_j$, the bisection schedule, the
quadrature tolerance) and the auditor logs are included in the
accompanying repository.

% ======================================================================
\section{Proof of the Main Theorem}\label{sec:proof}
% ======================================================================

\subsection{The Strict-Failure Witness}\label{subsec:proof-witness}

We close the master inequality using the anchors of
Table~\ref{tab:anchors}.  The strategy is to exhibit a rational
\emph{witness} $M_0=1292/1000$ for which $\Phi(M_0)<\tau$ holds
strictly; combined with the master inequality $\Phi(R(f))\ge\tau$,
this rules out $R(f)\le M_0$ by strict monotonicity of $\Phi$
(Lemma~\ref{lem:inversion}), forcing $R(f)>M_0$.  No analytic
optimization over $M$ is needed: the rational $M_0$ is fixed in
advance and the verification is purely arithmetic.
Set $\tau:=2/u+a$ and recall
$\Phi(M):=M+1+\sqrt{(M-1)(K_2-1)}$ from Lemma~\ref{lem:inversion}.
With the certified rationals from Table~\ref{tab:anchors},
\[
  \tau
  \;\ge\; \frac{2}{638/1000} + \frac{20925}{100000}
  \;=\; 3.134796\ldots + 0.20925
  \;=\; 3.344046\ldots
\]

\begin{proposition}[Strict failure at the rational witness]
\label{prop:fail}
With $K_2\le 47897/10000$ and $\tau=2/u+a\ge 2000/638+20925/10^{5}$,
\begin{equation}\label{eq:strict}
  \Phi(1292/1000) \;\le\; 3.34395
  \;<\; 3.344046\ldots \;\le\; \tau,
\end{equation}
with margin $\tau-\Phi(1292/1000)\;\ge\; 9.6\times 10^{-5}$.
\end{proposition}

\begin{proof}
At $M=1292/1000$, using the certified upper bound $K_2\le 47897/10000$,
\[
  (M-1)(K_2-1) \;\le\; \frac{292}{1000}\cdot\frac{37897}{10000}
  \;=\; \frac{11065924}{10^{7}}
  \;=\; 1.1065924,
\]
and since $\Phi$ is increasing in $K_2$, this upper bound transfers
to $\Phi(M)$.  The rational bound
$(105195/10^{5})^{2}=11065988025/10^{10}>11065924/10^{7}$
gives $\sqrt{(M-1)(K_2-1)}\le 105195/10^{5}=1.05195$, so
\[
  \Phi(1292/1000)
  \;\le\; \frac{1292}{1000} + 1 + \frac{105195}{10^{5}}
  \;=\; \frac{66879}{20000}
  \;=\; 3.34395.
\]
Combining with $\tau\ge 4267003/1276000=3.344046\ldots$ yields
\[
  \tau-\Phi(1292/1000)
  \;\ge\; \frac{4267003}{1276000} - \frac{66879}{20000}
  \;=\; \frac{307}{3190000}
  \;\ge\; 9.6\times 10^{-5}\;>\;0.
\]
The conclusion is independently re-certified by an arb cell-search
bisection at $256$-bit precision, which returns the sharper enclosure
$\Mcert\ge 1.29215$ when the script's tighter (non-rational) interval
bounds on $m_G$ and $S_1$ are used.
\end{proof}

\subsection{Assembly}\label{subsec:proof-assembly}

\begin{theorem}[Restatement of Theorem~\ref{thm:main}]
\label{thm:main-restated}
Let $f\in\mathcal{F}$ with $\|f*f\|_{\Linf}<\infty$.  Then
$R(f)\ge 1292/1000$, and hence $C_{1a}\ge 1292/1000$.
\end{theorem}

\begin{proof}
By homogeneity assume $\int f=1$, so $M:=R(f)=\|f*f\|_{\Linf}\ge 1$.
We assemble in five steps.
\begin{itemize}[leftmargin=2em]
\item[(1)] \emph{Admissibility.}  By Theorem~\ref{thm:admissibility},
$\Kms$ is Bochner-admissible at scale $\delta_1=138/1000$, with
period $u=638/1000$.
\item[(2)] \emph{Anchors.}  By Lemmas~\ref{lem:k1}--\ref{lem:a}, the
multiplier $G$ and the five rational bounds
$k_1\ge 9212/10000$, $K_2\le 47897/10000$, $S_1\le 29841/1000$,
$m_G\ge 998/1000$, and $a\ge 20925/100000$ hold.
\item[(3)] \emph{Master inequality.}  Theorem~\ref{thm:mv-master}
applied to $(\Kms,G)$ yields $\Phi(M)\ge \tau$.  Since $\Phi$ is
increasing in $K_2$, replacing the unknown $K_2$ by the certified
upper bound $47897/10000$ overestimates $\Phi$; the strict failure
$\Phi(M_0)<\tau$ at the witness $M_0=1292/1000$ therefore implies
$\Phi(R(f))>\Phi(M_0)$, which (by Lemma~\ref{lem:inversion}) forces
$R(f)>M_0$.
\item[(4)] \emph{Strict failure.}  By Proposition~\ref{prop:fail},
$\Phi(1292/1000)<\tau$.
\item[(5)] \emph{Conclusion.}  Strict monotonicity of $\Phi$ on
$[1,\infty)$ (Lemma~\ref{lem:inversion}) together with (3) and (4)
forces $M>1292/1000$.  Taking the infimum over $f\in\mathcal{F}$
gives $C_{1a}\ge 1292/1000=1.292$.\qedhere
\end{itemize}
\end{proof}

\subsection{Lean Formalization}\label{subsec:proof-lean}

The analytic chain above is mechanized in Lean~4 in the file
\texttt{lean/Sidon/MultiScale.lean}, which builds cleanly under
Mathlib~\cite{Mathlib,Lean} with zero \texttt{sorry} tactics and
produces the headline theorem.  The headline theorem reaches exactly
\emph{one} user axiom in its dependency closure:
\texttt{MV\_master\_inequality\_for\_extremiser}, the three-scale
extension of \cite[Lemma~3.1]{MV10} (whose proof is a verbatim repeat of
the single-scale Fourier reduction with each $J_0(\pi\delta\xi)^2$
replaced by $\sum_i\lambda_i J_0(\pi\delta_i\xi)^2$), stated with the
rational slacks $\texttt{K2UpperQ}=47897/10000$ and
$\texttt{gainLowerQ}=20925/100000$ substituted for $K_2$ and $a$
respectively.  This substitution is sound because the master inequality
is monotone in $K_2-1$ and $a$, and the slack rationals are true bounds
on the analytic functionals: those bounds (Lemmas~\ref{lem:k1}--\ref{lem:a})
are discharged externally by the \texttt{flint.arb} certifier in
\texttt{delsarte\_dual/grid\_bound\_alt\_kernel/}.

Six further statements are Lean \emph{theorems}: the algebraic
inversion \texttt{master\_inequality\_M\_lower} (corresponding to
Proposition~\ref{prop:fail}, proved by case analysis on $M\le 1$ versus
$M>1$ using \texttt{Real.sqrt} monotonicity) and the five
slack-soundness theorems \texttt{K\_two\_upper\_bound},
\texttt{k\_one\_lower\_bound}, \texttt{S\_one\_upper\_bound},
\texttt{min\_G\_lower\_bound}, \texttt{gain\_lower\_bound}, each a
one-line \texttt{norm\_num} verification that the rational slack is on
the correct side of the certifier-reported decimal recorded in
Lemmas~\ref{lem:k1}--\ref{lem:a}.

An explicit finiteness hypothesis $\|f*f\|_{\Linf}<\infty$ appears in
the Lean statement because of the \texttt{ENNReal.toReal} encoding;
this is harmless for the infimum.  The support hypothesis is stated
as $\texttt{Function.support}\,f\subseteq(-1/4,1/4)$, i.e.\
$f$ vanishes pointwise off $(-1/4,1/4)$; for any
$f\in\mathcal{F}$ in the sense of Section~\ref{subsec:intro-constant},
its representative obtained by zero-extending off the essential
support satisfies the pointwise condition and has the same
autoconvolution ratio, so the Lean bound transfers to the
$L^1$-equivalence-class formulation of $C_{1a}$ without loss.

% ======================================================================
\section*{Acknowledgments}
% ======================================================================

We thank M\'at\'e Matolcsi and Carlos Vinuesa, whose dual
framework~\cite{MV10} is the foundation of this argument; and
Alexander Cloninger and Stefan Steinerberger~\cite{CS17}, whose
announced bound of $1.2802$ motivated this line of work.  The
formalization rests on the Lean~4 theorem prover~\cite{Lean} and the
Mathlib library~\cite{Mathlib}; the certifier rests on
FLINT/Arb~\cite{Flint}.  We thank the authors and contributors of
each.


% ======================================================================
% References
% ======================================================================

\begin{thebibliography}{ABC10}

\bibitem[GGTW26]{AlphaEvolve25}
Bogdan Georgiev, Javier G\'omez-Serrano, Terence Tao, and Adam~Zsolt
Wagner, \emph{Mathematical exploration and discovery at scale},
preprint, 2026, arXiv:2511.02864,
\url{https://arxiv.org/abs/2511.02864}.

\bibitem[CS17]{CS17}
Alexander Cloninger and Stefan Steinerberger, \emph{On suprema of
autoconvolutions with an application to Sidon sets}, Proc. Amer. Math.
Soc.\ \textbf{145} (2017), no.~8, 3191--3200, arXiv:1403.7988.

\bibitem[ET41]{ET41}
Paul Erd\H{o}s and P\'al Tur\'an, \emph{On a problem of Sidon in additive
number theory, and on some related problems}, J. London Math. Soc.
\textbf{s1-16} (1941), no.~4, 212--215.

\bibitem[Joh17]{Flint}
Fredrik Johansson, \emph{Arb: efficient arbitrary-precision midpoint-radius
interval arithmetic}, IEEE Trans. Comput. \textbf{66} (2017), no.~8,
1281--1292.  See \url{http://arblib.org/}.

\bibitem[dMU21]{Lean}
Leonardo de Moura and Sebastian Ullrich, \emph{The Lean~4 theorem prover
and programming language}, in: Automated Deduction --- CADE~28, LNCS
\textbf{12699}, Springer, 2021, pp.~625--635.

\bibitem[MO09]{MO2009}
Greg Martin and Kevin O'Bryant, \emph{The supremum of autoconvolutions,
with applications to additive number theory}, Illinois J. Math.
\textbf{53} (2009), no.~1, 219--235, arXiv:0807.5121.

\bibitem[MV10]{MV10}
M\'at\'e Matolcsi and Carlos Vinuesa, \emph{Improved bounds on the supremum
of autoconvolutions}, J. Math. Anal. Appl. \textbf{372} (2010), no.~2,
439--447, arXiv:0907.1379.

\bibitem[The20]{Mathlib}
The mathlib community, \emph{The Lean mathematical library}, Proc.
9th ACM SIGPLAN Int. Conf. Certified Programs and Proofs, 2020,
pp.~367--381.

\bibitem[Wat44]{Watson}
George Neville Watson, \emph{A Treatise on the Theory of Bessel
Functions}, 2nd ed., Cambridge University Press, 1944.

\end{thebibliography}

\end{document}
